\medskip

\begin{enumerate}
\item  L'équation différentielle $y' = -4y + 80$ peut s'écrire $y' + 4y = 80$.
    
Il s'agit d'une équation différentielle du premier ordre de la forme :
\[y' + ay = b \text{ avec } a = 4 \text{ et } b = 80.\]

L'ensemble des solutions de $(E)$ est :
\begin{align*}
f(x) &= k\e^{-ax} + \dfrac{b}{a} \\
&= k\e^{-4x} + \dfrac{80}{4} \\
&= k\e^{-4x} + 20 \text{ où } k \in \mathbb{R}.
\end{align*}

\item Avec la condition initiale :
\begin{align*}
f(0) &= 100 \\
k\e^{-4\times0} + 20 &= 100 \\
k\e^0 &= 100 - 20 \\
k &= 80.
\end{align*}

La solution de $(E)$ est donc : $f(x) = 80\e^{-4x} + 20$.

\item Lorsque $x$ tend vers $+\infty$, $\e^{-4x}$ tend vers $0$. Par conséquent: $\displaystyle\lim_{x \to +\infty} f(x)=20$.
\end{enumerate}
